\begin{abstract}
Street-view perception models predict subjective attributes such as
safety at scale, yet they remain correlational: they do not reveal which
localised visual changes would plausibly shift human perception for a
specific street. We propose a lever-based interventional attribution
framework that reframes scene-level explainability as a structured search
over manipulable urban cues. Each candidate lever specifies a semantic
concept, spatial support, intervention direction, and constrained edit
template drawn from an editor-feasible intervention vocabulary. Candidates are realised via
prompt-conditioned diffusion editing and validated by a vision--language
critic for same-place preservation, locality, realism, and plausibility.
Across 50 scenes from five cities, 177 of 250 proposed candidates pass
the validity audit; all 50 scenes admit at least one valid
single-lever intervention and 48 admit multiple independently valid
levers. The strongest auxiliary perception shifts arise from
lane-marking repainting, crosswalk repainting, and localized greenery
addition, although valid realization does not guarantee positive
auxiliary movement. We position these results as evidence that bounded
interventional search is feasible at scene level, with human pairwise
evaluation remaining the planned endpoint.
\end{abstract}
